\chapter{Estado de la Cuestión}
\label{cap:estadoDeLaCuestion}

En las últimas décadas, el interés y la viabilidad de las interfaces cerebro-computador (Brain-Computer Interfaces, BCI) han aumentado de forma considerable gracias al avance de técnicas de aprendizaje automático y aprendizaje profundo aplicadas al análisis de bioseñales. Estos sistemas permiten establecer una comunicación entre el usuario y un dispositivo externo mediante la interpretación de la actividad cerebral, con aplicaciones en rehabilitación, asistencia a personas con discapacidad y nuevas formas de interacción hombre-máquina.

Para la adquisición de señales cerebrales pueden emplearse dispositivos invasivos y no invasivos. Los sistemas invasivos requieren procedimientos quirúrgicos para implantar sensores sobre o dentro del cerebro, como ocurre con la Electrocorticografía o los microelectrodos intracorticales. Por el contrario, los sistemas no invasivos registran la actividad cerebral desde el exterior sin necesidad de cirugía, destacando la Electroencefalografía (EEG) y la Magnetoencefalografía (MEG). Debido a su seguridad, menor coste y facilidad de uso, las técnicas no invasivas son las más empleadas en aplicaciones BCI, aunque presentan menor resolución espacial y mayor sensibilidad al ruido.

Existen diversos paradigmas BCI en función del mecanismo cerebral empleado. Entre los más relevantes se encuentran el paradigma P300, basado en la detección de un potencial evocado que aparece aproximadamente 300 ms después de un estímulo relevante, y el paradigma Steady-State Visual Evoked Potential (SSVEP), que aprovecha la respuesta cerebral sincronizada con estímulos visuales periódicos. Ambos paradigmas suelen ofrecer mayores tasas de precisión y menor tiempo de calibración, razón por la cual han sido ampliamente utilizados desde las primeras etapas del desarrollo de los sistemas BCI. Sin embargo, dependen de la presentación continua de estímulos visuales, auditivos o táctiles, lo que limita su naturalidad en según qué aplicaciones.

Por el contrario, el paradigma de imaginación motora (Motor Imagery, MI) permite generar comandos de forma voluntaria sin necesidad de estímulos externos. Aunque presenta habitualmente una mayor dificultad de clasificación y requiere un entrenamiento más prolongado por parte del usuario, resulta una alternativa especialmente atractiva para aplicaciones naturales, continuas o de rehabilitación, así como para su integración en videojuegos o entornos interactivos.

En el paradigma de imaginación motora, el usuario imagina la ejecución de un movimiento sin llegar a realizarlo físicamente, lo que provoca cambios detectables en la actividad cortical asociados a la planificación y representación motora. Una modalidad especialmente relevante es la Kinesthetic Motor Imagery, en la que el sujeto no solo representa visualmente la acción, sino que imagina la sensación interna del movimiento, incluyendo esfuerzo muscular, posición articular o desplazamiento del miembro. Estas variaciones se observan principalmente en regiones sensorimotoras del cerebro, localizadas alrededor del córtex motor primario, especialmente en áreas próximas a los electrodos C3, Cz y C4 del sistema internacional 10-20. Desde el punto de vista espectral, la información más relevante suele encontrarse en las bandas mu (8–13 Hz) y beta (13–30 Hz), donde se producen fenómenos de desincronización y resincronización relacionados con eventos motores. La detección de estos patrones permite diferenciar distintas tareas, como imaginar el movimiento de la mano izquierda, la mano derecha o los pies.

Dentro del ámbito de la clasificación EEG, uno de los principales desafíos actuales es la elevada variabilidad entre sujetos y entre sesiones, lo que dificulta la generalización de los modelos entrenados sobre múltiples usuarios. Esta variabilidad se debe tanto a diferencias anatómicas y fisiológicas entre individuos como a cambios en el estado del usuario a lo largo del tiempo. Para mitigar este problema, es habitual aplicar técnicas de preprocesado y normalización antes de la clasificación.

En cuanto al preprocesado, las señales EEG son especialmente sensibles a interferencias externas y artefactos fisiológicos, como los producidos por movimientos oculares o actividad muscular. Por ello, es habitual aplicar filtros paso banda para conservar únicamente las frecuencias de interés, en el caso de imaginación motora principalmente las bandas mu (8–13 Hz) y beta (13–30 Hz), así como técnicas de eliminación de artefactos.

Respecto a la normalización, se han propuesto diferentes estrategias para reducir las diferencias individuales en la actividad cerebral basal. Una de las más empleadas es la normalización por baseline, que consiste en restar o dividir la actividad registrada durante una ventana de reposo, tomada como referencia antes de cada tarea. Otra técnica relevante es el Euclidean Alignment, que alinea matemáticamente las señales de cada sujeto hacia una referencia común, permitiendo que el modelo entrenado con datos de otros usuarios se adapte mejor a uno nuevo sin necesidad de datos etiquetados adicionales, lo que resulta especialmente útil en entornos de uso real.

Uno de los modelos de referencia en clasificación EEG es EEGNet, una red neuronal convolucional diseñada para diferentes tareas de análisis EEG mediante un número reducido de parámetros. Más recientemente han surgido arquitecturas especializadas como MiRepNet, orientadas específicamente a la clasificación de imaginación motora, que buscan aprender representaciones más robustas a partir de grandes volúmenes de datos procedentes de distintos sujetos y conjuntos experimentales, permitiendo posteriormente adaptar el modelo mediante procesos de fine-tuning con una cantidad reducida de datos del usuario final.

En el estado de la cuestión es donde aparecen gran parte de las referencias bibliográficas del trabajo. Una de las formas más cómodas de gestionar la bibliografía en {\LaTeX} es utilizando \textbf{bibtex}. Las entradas bibliográficas deben estar en un fichero con extensión \textit{.bib} (con esta plantilla se proporciona el fichero biblio.bib, donde están las entradas referenciadas más abajo). Cada entrada bibliográfica tiene una clave que permite referenciarla desde cualquier parte del texto con los siguiente comandos:

\begin{itemize}
\item Referencia bibliografica con cite: \cite{ldesc2e}
\item Referencia bibliográfica con citep: \citep{notsoshort}
\item Referencia bibliográfica con citet: \citet{latexAPrimer}
\end{itemize}

Es posible citar más de una fuente, como por ejemplo \citep{latexCompanion,LaTeXLamport,texKnuth}

Después, \LaTeX se ocupa de rellenar la sección de bibliografía con las entradas \textbf{que hayan sido citadas} (es decir, no con todas las entradas que hay en el .bib, sino sólo con aquellas que se hayan citado en alguna parte del texto).

Bibtex es un programa separado de latex, pdflatex o cualquier otra cosa que se use para compilar los .tex, de manera que para que se rellene correctamente la sección de bibliografía es necesario compilar primero el trabajo (a veces es necesario compilarlo dos veces), compilar después con bibtex, y volver a compilar otra vez el trabajo (de nuevo, puede ser necesario compilarlo dos veces). 
